\par\Needspace{6\baselineskip}\medskip
\begin{center}\hypertarget{native-178-article131}{}\pdfbookmark[2]{131.}{article-131}131.\end{center}
In Artikel 114 wurde gezeigt, dass sich die Abstände des Himmelskörpers von der Erde mit sehr
einfachen Formeln bestimmen lassen, wenn die Verhältnisse zwischen den dort und in Artikel 128 mit
$n,n',n''$ bezeichneten Größen bekannt sind. Nähme man also für $x,y$ die Quotienten $\dfrac{n}{n'},\ \dfrac{n''}{n'}$, so würden sich,
sofern die heliozentrische Bewegung zwischen den Beobachtungen nicht sehr groß ist, sogleich die
Näherungswerte $\dfrac{\theta}{\theta'},\ \dfrac{\theta''}{\theta'}$ anbieten, wobei die Zeichen $\theta,\theta',\theta''$ dieselbe Bedeutung wie in Artikel 128 haben.
Damit scheint sich eine naheliegende Lösung unserer Aufgabe zu ergeben, wenn man aus $x$ und $y$
zwei Abstände von der Erde ermittelt und dann nach einem der fünf Verfahren aus den Artikeln 124–128
vorgeht. Behält man auch für die Zeichen $\eta,\eta''$ die Bedeutung aus Artikel 128 bei und bezeichnet man
entsprechend mit $\eta'$ den Quotienten aus dem Flächeninhalt des von zwei Radiusvektoren
eingeschlossenen Sektors und dem Flächeninhalt des Dreiecks zwischen denselben Radiusvektoren, so
erhält man tatsächlich
\[
\frac{n}{n'}=\frac{\theta}{\theta'}\cdot\frac{\eta'}{\eta}\mathpunct{\raisebox{.45ex}{,}}\qquad
\frac{n''}{n'}=\frac{\theta''}{\theta'}\cdot\frac{\eta'}{\eta''}\mathpunct{\raisebox{.45ex}{,}}
\]
und betrachtet man $n,n',n''$ als kleine Größen erster Ordnung, so ist leicht zu erkennen, dass $\eta-1,\eta'-1,\eta''-1$ im
Allgemeinen Größen zweiter Ordnung sind. Die Näherungswerte für $x,y$, nämlich $\dfrac{\theta}{\theta'},\ \dfrac{\theta''}{\theta'}$, unterscheiden
sich daher von den wahren Werten nur um Größen zweiter Ordnung. Bei näherer Betrachtung erweist sich
dieses Verfahren dennoch als völlig ungeeignet. Den Grund dafür wollen wir kurz erläutern. Man sieht
leicht, dass die Größe $(0.1.2)$, mit der die Abstände in den Formeln 9, 10 und 11 aus Artikel 114
multipliziert sind, mindestens von dritter Ordnung ist. In Gleichung 9 etwa sind die Größen $(\mathrm{O}.1.2)$,
$(\mathrm{I}.1.2)$ und $(\mathrm{II}.1.2)$ von erster Ordnung. Daraus folgt leicht, dass ein Fehler zweiter Ordnung in den Werten
von $\dfrac{n}{n'},\ \dfrac{n''}{n'}$ in den Werten der Abstände einen Fehler der Ordnung $0$ hervorruft. In der üblichen
Ausdrucksweise wären die so erhaltenen Abstände also selbst dann mit einem endlichen Fehler
behaftet, wenn die Zeitintervalle unendlich klein sind. Folglich ließen sich weder diese Abstände
noch die übrigen daraus abzuleitenden Größen auch nur als Näherungswerte ansehen, und das Verfahren
würde die zweite Bedingung des vorigen Artikels verletzen.
